\documentclass[12pt,a4paper]{article}
\usepackage[left=1.5cm,right=1.5cm,top=2cm,bottom=2.5cm]{geometry}
\usepackage{multicol}
\usepackage{lipsum} % لتوليد نص تجريبي
% 1. استدعاء حزمة الزخارف المعدلة
\usepackage{na-fancyborders}
\setlength{\columnsep}{2cm}
\usepackage{polyglossia}
\setmainlanguage[numerals=maghrib]{arabic}
\setotherlanguage{french}

% 3. إعدادات الخطوط (تأكد من وجود الخطوط على جهازك)
\newfontfamily\arabicfont[Script=Arabic,Scale=1.2]{Amiri}
\newfontfamily\arabicfontsf[Script=Arabic,Scale=1.2]{Arial}

% 4. إعدادات الهوامش (Geometry) لتناسب الإطارات


\begin{document}

% تفعيل الإطار الأول
\activateBGdefault
\section*{الصفحة الأولى}
هنا تكتب النص العربي الخاص بك. هذا الإطار سيطبق على هذه الصفحة فقط.
\newpage
% تنظيف الصفحة من الإطار السابق وتفعيل الإطار الثاني
\ClearShipoutPicture
\activateBGone
\section*{الصفحة الثانية}
هنا نص الصفحة الثانية بإطار مختلف تماماً.
\newpage
% تنظيف الصفحة من الإطار السابق وتفعيل الإطار الثاني
\ClearShipoutPicture
\activateBGtwo
\section*{الصفحة الثانية}
هنا نص الصفحة الثانية بإطار مختلف تماماً.

\newpage
\ClearShipoutPicture
\activateBGfour
\section*{الصفحة الثالثة}
هذا هو الإطار الذي يحتوي على ترقيم الصفحة في الأسفل بشكل فني.
\newpage
\ClearShipoutPicture       % مسح أي إطارات سابقة
\activateBGthree           % تفعيل الزخارف الوسطية (بدون إطار خارجي)

% --- بدء نظام العمودين ---
\begin{multicols*}{2}
\lipsum[1-4]
\end{multicols*}
\end{document}